\documentclass[10pt,a4paper]{amsart}
\usepackage{amsmath,amssymb,amsthm,graphicx}
\usepackage{verbatim}
\newenvironment{reference}{\comment}{\endcomment}
%\newenvironment{reference}{}{}
\newenvironment{slogan}{\comment}{\endcomment}
\newenvironment{history}{\comment}{\endcomment}

% For commutative diagrams we use Xy-pic
\usepackage[all]{xy}

% We use 2cell for 2-commutative diagrams.
\xyoption{2cell}
% Theorem environments.
%
\theoremstyle{plain}
\newtheorem{theorem}[subsection]{Theorem}
\newtheorem{proposition}[subsection]{Proposition}
\newtheorem{lemma}[subsection]{Lemma}

\theoremstyle{definition}
\newtheorem{definition}[subsection]{Definition}
\newtheorem{example}[subsection]{Example}
\newtheorem{exercise}[subsection]{Exercise}
\newtheorem{situation}[subsection]{Situation}

\theoremstyle{remark}
\newtheorem{remark}[subsection]{Remark}
\newtheorem{remarks}[subsection]{Remarks}

\numberwithin{equation}{subsection}

% Macros
%
\def\lim{\mathop{\mathrm{lim}}\nolimits}
\def\colim{\mathop{\mathrm{colim}}\nolimits}
\def\Spec{\mathop{\mathrm{Spec}}}
\def\Hom{\mathop{\mathrm{Hom}}\nolimits}
\def\Ext{\mathop{\mathrm{Ext}}\nolimits}
\def\SheafHom{\mathop{\mathcal{H}\!\mathit{om}}\nolimits}
\def\SheafExt{\mathop{\mathcal{E}\!\mathit{xt}}\nolimits}
\def\Sch{\mathit{Sch}}
\def\Mor{\mathop{\mathrm{Mor}}\nolimits}
\def\Ob{\mathop{\mathrm{Ob}}\nolimits}
\def\Sh{\mathop{\mathit{Sh}}\nolimits}
\def\NL{\mathop{N\!L}\nolimits}
\def\CH{\mathop{\mathrm{CH}}\nolimits}
\def\proetale{{pro\text{-}\acute{e}tale}}
\def\etale{{\acute{e}tale}}
\def\QCoh{\mathit{QCoh}}
\def\Ker{\mathop{\mathrm{Ker}}}
\def\Im{\mathop{\mathrm{Im}}}
\def\Coker{\mathop{\mathrm{Coker}}}
\def\Coim{\mathop{\mathrm{Coim}}}

% Boxtimes
%
\DeclareMathSymbol{\boxtimes}{\mathbin}{AMSa}{"02}

%
% Macros for moduli stacks/spaces
%
\def\QCohstack{\mathcal{QC}\!\mathit{oh}}
\def\Cohstack{\mathcal{C}\!\mathit{oh}}
\def\Spacesstack{\mathcal{S}\!\mathit{paces}}
\def\Quotfunctor{\mathrm{Quot}}
\def\Hilbfunctor{\mathrm{Hilb}}
\def\Curvesstack{\mathcal{C}\!\mathit{urves}}
\def\Polarizedstack{\mathcal{P}\!\mathit{olarized}}
\def\Complexesstack{\mathcal{C}\!\mathit{omplexes}}
% \Pic is the operator that assigns to X its picard group, usage \Pic(X)
% \Picardstack_{X/B} denotes the Picard stack of X over B
% \Picardfunctor_{X/B} denotes the Picard functor of X over B
\def\Pic{\mathop{\mathrm{Pic}}\nolimits}
\def\Picardstack{\mathcal{P}\!\mathit{ic}}
\def\Picardfunctor{\mathrm{Pic}}
\def\Deformationcategory{\mathcal{D}\!\mathit{ef}}
\usepackage{cleveref}
\crefname{originallemma}{original lemma}{original lemmas}
\Crefname{originallemma}{Original lemma}{Original lemmas}
\makeatletter
\newlabel{algebra-lemma-exists-dvr}{{00PH}{source}{}{sourcepointer-1}{}}
\newlabel{algebra-lemma-exists-dvr@cref}{{[originallemma][0][]00PH}{[1][0][]source}{}{}{}}
\makeatother
\crefname{originallemma}{original lemma}{original lemmas}
\Crefname{originallemma}{Original lemma}{Original lemmas}
\makeatletter
\newlabel{formal-spaces-lemma-morphism-between-formal-spectra}{{0AN0}{source}{}{sourcepointer-2}{}}
\newlabel{formal-spaces-lemma-morphism-between-formal-spectra@cref}{{[originallemma][0][]0AN0}{[1][0][]source}{}{}{}}
\makeatother
\crefname{originallemma}{original lemma}{original lemmas}
\Crefname{originallemma}{Original lemma}{Original lemmas}
\makeatletter
\newlabel{properties-lemma-complement-closed-point-Jacobson}{{02IM}{source}{}{sourcepointer-3}{}}
\newlabel{properties-lemma-complement-closed-point-Jacobson@cref}{{[originallemma][0][]02IM}{[1][0][]source}{}{}{}}
\makeatother
\crefname{originalreference}{original reference}{original references}
\Crefname{originalreference}{Original reference}{Original references}
\makeatletter
\newlabel{remark-base-change}{{06FQ}{source}{}{sourcepointer-4}{}}
\newlabel{remark-base-change@cref}{{[originalreference][0][]06FQ}{[1][0][]source}{}{}{}}
\makeatother
\title{Flat reductions with faithful flatness in degree one give rig-surjective formal spectra}
\author{The Stacks Project Authors}
\date{stacks-project:restricted.tex:local-snapshot-20261012; GNU FDL 1.2 or later}
\begin{document}
\maketitle
\noindent\textit{Editorial note.} All delivered mathematical text is the author's own, copied from the pinned Stacks Project chapter file restricted.tex; no statement or proof step was rewritten, reordered or removed.\par\medskip
\noindent\textit{Editorial note.} The counted proof is the author's complete argument: the source proves the lemma in full, constructing the lift via the completed base change and a one-dimensional quotient embedded into a DVR\par\medskip
\noindent\textit{Editorial note.} The label lemma-faithfully-flat-rig-surjective names the counted statement, and the proof environment immediately following it in the source is delivered verbatim as its proof.\par\medskip
\noindent\textit{Editorial note.} The counted statement is carried by the author's own lemma environment in the chapter source.\par\medskip
\noindent\textit{Allowed standard prerequisites (editorial).}\begin{itemize}
\item formal spectra of adic Noetherian rings and morphisms given by continuous ring maps
\item rig-surjectivity tested by lifts along adic maps to complete discrete valuation rings
\item the completed base change C = B completed tensor\_A R and its reduction properties
\item going-down for flat maps and the Jacobson complement-of-a-closed-point lemma
\item existence of a one-dimensional local domain quotient embedding into a discrete valuation ring
\end{itemize}\medskip
\ifdefined\hypertarget\hypertarget{sourcepointer-1}{}\fi
\noindent\textit{Source reference.} Original lemma 00PH: sources/stacks-chow/tags.txt, source lines 931--931. Prerequisite pointer invoked at proof line(s) 5451; it names a standard fact of a sibling chapter that the counted argument uses, and every step of the counted proof built on it is retained verbatim in the delivered body.\par\medskip
\ifdefined\hypertarget\hypertarget{sourcepointer-2}{}\fi
\noindent\textit{Source reference.} Original lemma 0AN0: sources/stacks-chow/tags.txt, source lines 13423--13423. Prerequisite pointer invoked at proof line(s) 5429; it names a standard fact of a sibling chapter that the counted argument uses, and every step of the counted proof built on it is retained verbatim in the delivered body.\par\medskip
\ifdefined\hypertarget\hypertarget{sourcepointer-3}{}\fi
\noindent\textit{Source reference.} Original lemma 02IM: sources/stacks-chow/tags.txt, source lines 3305--3305. Prerequisite pointer invoked at proof line(s) 5448; it names a standard fact of a sibling chapter that the counted argument uses, and every step of the counted proof built on it is retained verbatim in the delivered body.\par\medskip
\ifdefined\hypertarget\hypertarget{sourcepointer-4}{}\fi
\noindent\textit{Source reference.} Original reference 06FQ: sources/stacks-restricted/pinned-original-source.tex, source lines 426--426. Prerequisite pointer invoked at proof line(s) 5436; it names a standard fact of the same chapter that the counted argument uses, and every step of the counted proof built on it is retained verbatim in the delivered body.\par\medskip
\begin{lemma}
\label{lemma-faithfully-flat-rig-surjective}
Let $A$ be a Noetherian ring complete with respect to an ideal $I$.
Let $B$ be an $I$-adically complete $A$-algebra.
If $A/I^n \to B/I^nB$ is of finite type and flat for all $n$ and
faithfully flat for $n = 1$, then $\text{Spf}(B) \to \text{Spf}(A)$
is rig-surjective.
\end{lemma}

\begin{proof}
We will use without further mention that morphisms between formal spectra
are given by continuous maps between the corresponding topological rings, see
Formal Spaces, Lemma \ref{formal-spaces-lemma-morphism-between-formal-spectra}.
Let $\varphi : A \to R$ be a continuous map into a complete discrete
valuation ring $A$. This implies that $\varphi(I) \subset \mathfrak m_R$.
On the other hand, since we only need to produce the lift
$\varphi' : B' \to R'$ in the case that $\varphi$ corresponds to an adic
morphism, we may assume that $\varphi(I) \not = 0$. Thus we may consider
the base change $C = B \widehat{\otimes}_A R$, see
Remark \ref{remark-base-change} for example.
Then $C$ is an $\mathfrak m_R$-adically complete $R$-algebra
such that $C/\mathfrak m_R^n C$ is of finite type and flat over
$R/\mathfrak m_R^n$ and such that $C/\mathfrak m_R C$ is nonzero.
Pick any maximal ideal $\mathfrak m \subset C$ lying over
$\mathfrak m_R$. By flatness (which implies going down) we see that
$\Spec(C_\mathfrak m) \setminus V(\mathfrak m_R C_\mathfrak m)$
is a nonempty open. Hence
We can pick a prime $\mathfrak q \subset \mathfrak m$
such that $\mathfrak q$ defines a closed point of
$\Spec(C_\mathfrak m) \setminus \{\mathfrak m\}$ and such that
$\mathfrak q \not \in V(IC_\mathfrak m)$, see
Properties, Lemma \ref{properties-lemma-complement-closed-point-Jacobson}.
Then $C/\mathfrak q$ is a dimension $1$-local domain and we can find
$C/\mathfrak q \subset R'$ with $R'$ a discrete valuation ring
(Algebra, Lemma \ref{algebra-lemma-exists-dvr}).
By construction $\mathfrak m_R R' \subset \mathfrak m_{R'}$
and we see that $C \to R'$ extends to a continuous map
$C \to (R')^\wedge$ (in fact we can pick $R'$ such that
$R' = (R')^\wedge$ in our current situation but we do not need this).
Since the completion of a discrete valuation ring is a discrete
valuation ring, we see that the assumption gives a commutative
diagram of rings
$$
\xymatrix{
(R')^\wedge & C \ar[l] & B \ar[l] \\
R \ar[u] & R \ar[l] \ar[u] & A \ar[l] \ar[u]
}
$$
which gives the desired lift.
\end{proof}

\end{document}
